\chapter{Conception et implémentation de la solution}
\label{chap:implementation}

\section{Démarche d'ingénierie}

Mon travail a suivi une démarche itérative fondée sur l'observation des
exécutions plutôt que sur l'ajout isolé de fonctionnalités. Pour chaque
évolution, j'ai d'abord identifié un défaut reproductible dans une trace, un
artefact de phase ou une campagne de benchmark. J'ai ensuite localisé la
frontière responsable --- modèle, contrat d'outil, orchestration ou évaluateur
--- avant d'implémenter un correctif et d'ajouter un test de non-régression.
Cette méthode était particulièrement importante pour évaluer l’agent génératif intégré à LANCE: une sortie conforme au format attendu peut néanmoins être techniquement erronée. Inversement, un échec observé ne provient pas nécessairement du modèle, mais peut résulter d’un contexte d’entrée tronqué ou d’un contrat d’interface ambigu.


Les tests ont porté sur les frontières les plus risquées : validité des données
d'entraînement, compatibilité de l'environnement, sélection de l'expert,
respect des outils exposés, reprise après rejet d'un livrable, cohérence des
artefacts entre phases et calcul des scores. Le but n'était pas seulement de
faire terminer une campagne, mais de rendre explicite la raison pour laquelle
elle pouvait être considérée comme complète.

\section{Architecture générale}

LANCE sépare les décisions probabilistes des contrôles déterministes. La classe
\texttt{\footnotesize Pipeline} orchestre six phases et associe à chacune un prompt, un
livrable attendu et un ensemble d'outils autorisés. Le modèle propose une
analyse ou une action ; l'orchestrateur exécute l'outil, journalise son résultat
et décide si le contrat de la phase est satisfait. Les artefacts structurés
forment ainsi une mémoire interphase plus fiable qu'une conversation libre.

Autour du pipeline, FastAPI expose le lancement et le suivi des campagnes. Un
flux \gls{sse} alimente l'interface, qui affiche la topologie, les livrables et
les métriques. Le benchmark déploie les scénarios au moyen d'Ansible sur
Proxmox, tandis qu'un évaluateur distinct compare les artefacts à la vérité terrain.
Enfin, un serveur \gls{hmoe} compatible avec l'\gls{api} OpenAI fournit au pipeline un point
d'accès unique à plusieurs adaptateurs spécialisés.

La \cref{fig:architecture-lance} présente les principales frontières du système :
le déploiement fournit la cible au pipeline, tandis que la vérité terrain reste
du côté de l'évaluateur. Le pipeline conserve seul la responsabilité d'autoriser
les outils, d'exécuter les actions et de produire les artefacts évalués.

\begin{figure}[htbp]
  \centering
  \resizebox{\textwidth}{!}{\usetikzlibrary{shapes, arrows.meta, positioning, fit}

\begin{tikzpicture}[
  font=\sffamily\small,
  >={Latex[length=3mm, width=2.2mm]},
  flow/.style={->, thick, draw=black!75},
  feedback/.style={->, thick, dashed, draw=violet!80!black},
  component/.style={draw=blue!75!black, fill=blue!8, rounded corners=3pt,
    align=center, minimum height=11mm, text width=30mm, inner sep=5pt, semithick},
  phase/.style={draw=cyan!75!black, fill=cyan!8, rounded corners=2pt,
    align=center, minimum height=10mm, text width=22mm, inner sep=4pt, font=\sffamily\footnotesize, semithick},
  external/.style={draw=orange!80!black, fill=orange!10, rounded corners=3pt,
    align=center, minimum height=11mm, text width=30mm, inner sep=5pt, semithick},
  artifact/.style={draw=green!60!black, fill=green!10, rounded corners=3pt,
    align=center, minimum height=11mm, text width=32mm, inner sep=5pt, semithick},
  group/.style={draw=black!50, dashed, rounded corners=4pt, inner sep=8pt, semithick},
  lbl/.style={font=\sffamily\scriptsize, fill=white, inner sep=1.5pt} % Style pour texte sur les flèches
]

  % --- Briques externes (Gauche) ---
  \node[external] (catalogue) {Catalogue de scénarios\\et vérité terrain};
  \node[external, below=12mm of catalogue] (deploy) {Ansible + Proxmox\\déploiement reproductible};
  \node[external, below=12mm of deploy] (target) {Réseau cible\\IoT/OT};

  % --- Interface & API ---
  \node[component, right=22mm of catalogue] (ui) {Tableau de bord\\opérateur};
  \node[component, below=12mm of ui] (api) {API FastAPI\\suivi par SSE};

  % --- Pipeline LANCE ---
  \node[phase, right=18mm of api, yshift=24mm] (p1) {1. Analyse\\du graphe};
  \node[phase, right=6mm of p1] (p2) {2. Recon-\\naissance};
  \node[phase, right=6mm of p2] (p3) {3. Vulnéra-\\bilités};
  \node[phase, below=10mm of p3] (p4) {4. Vérifi-\\cation};
  \node[phase, left=6mm of p4] (p5) {5. Intrusion\\contrôlée};
  \node[phase, left=6mm of p5] (p6) {6. Rapport\\final};
  
  \node[group, fit=(p1)(p2)(p3)(p4)(p5)(p6),
    label={[font=\sffamily\small\bfseries, yshift=2pt]above:Pipeline LANCE}] (pipeline) {};

  % --- Modèle & Outils ---
  \node[component, right=18mm of p3] (model) {Fournisseur LLM\\HMoE local ou API distante};
  \node[component, below=12mm of model] (tools) {Registre d'outils\\Nmap, HTTP, MQTT,\\Modbus, SSH, \ldots};

  % --- Artefacts & Évaluation ---
  \node[artifact, below=18mm of p6, xshift=15mm] (artifacts) {Artefacts structurés\\et traces d'appels};
  \node[artifact, right=18mm of artifacts] (harness) {Harness strict-v3\\comparaison à la vérité terrain};
  \node[artifact, right=18mm of harness] (metrics) {Scores, diagnostics\\et rapports de campagne};

  % --- Liens et Flux ---
  \draw[flow] (catalogue) -- (deploy);
  \draw[flow] (deploy) -- (target);
  \draw[flow] (ui) -- (api);
  \draw[flow] (api.east) -- (pipeline.west);
  
  \draw[flow] (p1) -- (p2);
  \draw[flow] (p2) -- (p3);
  \draw[flow] (p3) -- (p4);
  \draw[flow] (p4) -- (p5);
  \draw[flow] (p5) -- (p6);
  
  \draw[flow, bend left=15] (pipeline.east) to node[lbl, above=1pt]{requête} (model.west);
  \draw[feedback, bend left=15] (model.west) to node[lbl, below=1pt]{réponse / appel} (pipeline.east);
  
  \draw[flow] (pipeline.south east) -- (tools.north west);
  \draw[flow] (tools.west) -- node[lbl, below]{actions bornées} (target.east);
  \draw[feedback] (target.north east) -- node[lbl, above]{observations} (pipeline.south west);
  
  \draw[flow] (pipeline.south) -- (artifacts.north);
  \draw[flow] (artifacts) -- (harness);
  \draw[flow] (harness) -- (metrics);
  
  \draw[flow] (catalogue.south east) |- node[lbl, pos=0.8, above]{référence scellée} (harness.north);
  \draw[feedback] (metrics.north west) |- (ui.south);

\end{tikzpicture}
}
  \caption{Architecture fonctionnelle de LANCE et de son environnement d'évaluation}
  \label{fig:architecture-lance}
\end{figure}


\section{Mission principale : construction d'un HMoE agentique}

\subsection{Principe de spécialisation}

La mission principale de mon stage a consisté à construire une chaîne HMoE
agentique, depuis les traces d'exécution jusqu'au déploiement des experts dans
LANCE. Ici, le terme HMoE ne désigne pas une couche de mélange d'experts apprise
à l'intérieur du réseau. Le système partage un modèle de base, charge au
démarrage quatre adaptateurs \gls{peft} spécialisés, puis active celui qui
correspond à chaque requête : \texttt{\footnotesize secretary}, \texttt{\footnotesize recon},
\texttt{\footnotesize vuln} et \texttt{\footnotesize exploit}. Cette organisation reprend
l'idée de spécialisation des mélanges d'experts tout en l'adaptant aux
contraintes matérielles du projet \cite{shazeer2017moe}.

Le secrétaire traite l'analyse initiale du graphe et la synthèse finale. L'expert
de reconnaissance couvre la phase~2 ; l'expert de vulnérabilités intervient en
phase~3 ; l'expert d'exploitation couvre les phases~4 et~5. Cette répartition
évite d'entraîner chaque adaptateur sur l'intégralité des outils et des formats.
Elle réduit aussi l'ambiguïté : une trace d'exploitation et une trace de rapport
ne demandent ni les mêmes actions ni le même critère de terminaison.

Le choix d'adaptateurs plutôt que de modèles indépendants répond à une
contrainte de mémoire GPU. Une seule instance quantifiée du modèle de base est
chargée, puis l'adaptateur correspondant est sélectionné pour chaque requête. Le gain de
mémoire a une contrepartie : le changement d'adaptateur modifie un état partagé.
J'ai donc sérialisé la génération avec un verrou. Les requêtes peuvent être
reçues concurremment par l'API, mais deux générations ne changent jamais
l'adaptateur au même instant.

La \cref{fig:architecture-hmoe} détaille ce routage. La sélection porte sur une
requête complète et précède la génération : les quatre experts partagent donc
les poids quantifiés du même modèle, mais appliquent chacun une correction
QLoRA propre à leur rôle.

\begin{figure}[htbp]
  \centering
  \resizebox{0.98\textwidth}{!}{\begin{tikzpicture}[
  font=\sffamily\small,
  >={Latex[length=2.8mm,width=2mm]},
  flow/.style={->,thick,draw=black!75},
  link/.style={thick,draw=black!65},
  return/.style={->,thick,dashed,draw=violet!80!black},
  block/.style={draw=blue!70!black,fill=blue!8,rounded corners=3pt,
    align=center,minimum height=11mm,text width=31mm,inner sep=5pt,semithick},
  expert/.style={draw=green!55!black,fill=green!9,rounded corners=3pt,
    align=center,minimum height=10mm,text width=24mm,inner sep=4pt,semithick},
  guard/.style={draw=orange!80!black,fill=orange!9,rounded corners=3pt,
    align=center,minimum height=11mm,text width=94mm,inner sep=5pt,semithick},
  group/.style={draw=black!50,dashed,rounded corners=4pt,inner sep=7pt,semithick},
  lbl/.style={font=\sffamily\scriptsize,fill=white,inner sep=1.5pt}
]
  % Une requête entière est routée avant toute génération
  \node[block] (pipeline) {Pipeline LANCE\\requête OpenAI-compatible};
  \node[block,right=20mm of pipeline] (router)
    {Routeur déterministe\\phase et rôle\\jamais par jeton};
  \coordinate (routingmid) at ($(pipeline)!0.5!(router)$);
  \coordinate (split) at ($(routingmid)+(0,-18mm)$);

  % Quatre adaptateurs spécialisés clairement séparés
  \node[expert,below=20mm of routingmid,xshift=-54mm] (secretary)
    {\texttt{secretary}\\phases 1 et 6};
  \node[expert,right=8mm of secretary] (recon)
    {\texttt{recon}\\phase 2};
  \node[expert,right=8mm of recon] (vuln)
    {\texttt{vuln}\\phase 3};
  \node[expert,right=8mm of vuln] (exploit)
    {\texttt{exploit}\\phases 4 et 5};
  \node[group,fit=(secretary)(recon)(vuln)(exploit),
    label={[font=\sffamily\small\bfseries]above:Adaptateurs QLoRA spécialisés}] (adapters) {};

  % Socle commun et exécution protégée
  \node[block,below=18mm of adapters] (base)
    {Qwen2.5-3B-Instruct\\base quantifiée 4 bits\\partagée et gelée};
  \node[block,below=14mm of base] (generation)
    {Adaptateur actif\\+ génération sérialisée};
  \node[guard,below=14mm of generation] (guards)
    {Garde-fous : budget de contexte et de génération, filtrage des outils\\
     autorisés et reprise contrôlée après un rejet};

  % Routage de la requête
  \draw[flow] (pipeline) -- node[above,font=\scriptsize]{requête complète} (router);
  \draw[flow] (router.south) |- (split);
  \draw[flow] (split) -| (secretary.north);
  \draw[flow] (split) -| (recon.north);
  \draw[flow] (split) -| (vuln.north);
  \draw[flow] (split) -| (exploit.north);

  % Convergence vers l'unique modèle de base
  \coordinate (merge) at ($(base.north)+(0,8mm)$);
  \draw[link] (secretary.south) |- (merge);
  \draw[link] (recon.south) |- (merge);
  \draw[link] (vuln.south) |- (merge);
  \draw[link] (exploit.south) |- (merge);
  \draw[flow] (merge) -- (base.north);
  \draw[flow] (base) -- (generation);
  \draw[return] (guards.north) -- node[lbl,right]{contraintes} (generation.south);
  \draw[return] (generation.west) -- ++(-62mm,0) |-
    node[lbl,pos=0.25,below]{texte ou appel d'outil} (pipeline.south);
\end{tikzpicture}
}
  \caption{Architecture du HMoE agentique et routage déterministe par phase}
  \label{fig:architecture-hmoe}
\end{figure}

\subsection{Construction des données d'apprentissage}

La qualité du système dépend d'abord de la correspondance entre les données et
la boucle agentique réellement exécutée. J'ai transformé les traces en exemples
de conversation contenant les messages, les appels d'outils, leurs résultats et
les schémas d'outils disponibles. Les exemples sont routés vers l'expert associé
à la phase qui a produit le livrable. Cette conservation des appels est
essentielle : un corpus qui ne garderait que la réponse finale apprendrait à
rédiger, mais pas à décider quand observer le réseau ou quand sauvegarder un
artefact.

La préparation cible Qwen2.5-3B-Instruct et applique le gabarit de conversation
utilisé à l'inférence. Pour chaque exemple, seuls les outils effectivement
appelés et quelques distracteurs sont conservés. Les distracteurs obligent le
modèle à sélectionner un outil pertinent sans gonfler inutilement le contexte.
Les conversations longues sont découpées à la frontière des tours assistant ;
les couples appel/résultat restent liés. La phase~5 constitue une exception : sa
chaîne de preuves est gardée atomique afin qu'une conclusion d'intrusion ne soit
jamais séparée de l'action qui la soutient.

Lorsque la fenêtre reste trop longue, la préparation réduit d'abord les sorties
d'outils volumineuses, puis les arguments textuels et enfin le milieu des
messages système. Le préfixe contenant le rôle et le suffixe décrivant le
critère de terminaison sont préservés. Chaque ligne indique dans ses métadonnées
le modèle cible, la fenêtre source et l'éventuelle troncature. Les limites sont
alignées sur l'usage futur : 6\,144 tokens pour le secrétaire, la reconnaissance
et l'exploitation, et 4\,096 pour l'analyse de vulnérabilités.

J'ai également intégré une boucle de correction supervisée. Le harness peut
extraire des candidats correspondant à un faux négatif, un faux positif, une
sévérité incorrecte ou un livrable incohérent. Ces candidats restent en attente
jusqu'à une revue humaine ; seuls ceux explicitement acceptés sont convertis en
exemples SFT. Ils sont attribués à la première phase disposant d'un signal causal
observable. Une cible absente de la reconnaissance corrige ainsi l'expert
\texttt{\footnotesize recon}, et non automatiquement l'expert \texttt{\footnotesize vuln}. Ce choix évite
d'enseigner à un composant aval à compenser silencieusement une erreur amont.

\subsection{Préflight et reproductibilité}

Avant tout entraînement, un préflight vérifie les conditions sans charger les
poids ni lancer d'optimisation. Il contrôle la structure YAML, la présence du
gabarit et de ses masques de génération, puis parcourt les JSONL. Chaque ligne
doit contenir des messages, au moins une réponse assistant et des métadonnées
cohérentes avec l'expert et le modèle préparé. Toute référence au split
\texttt{\footnotesize eval-sealed} est refusée afin d'empêcher une contamination des données
d'évaluation.

Le préflight contrôle aussi les versions de PyTorch, Transformers, TRL, PEFT,
bitsandbytes, Datasets et Accelerate. Il vérifie que la signature de
\texttt{\footnotesize SFTConfig} expose les paramètres requis par le script, que CUDA est
disponible, que le GPU accepte BF16 et, en mode strict, qu'aucun autre processus
de calcul ne l'occupe. Il produit un rapport JSON comprenant les tailles et
empreintes des jeux de données. Lors d'une promotion, le même outil exige pour
chaque expert un fichier de configuration, les poids finaux et le tokenizer,
puis confirme que l'adaptateur référence le modèle de base prévu.


\subsection{Entraînement QLoRA}

J'ai retenu QLoRA pour adapter le modèle avec une empreinte mémoire compatible
avec la machine disponible \cite{dettmers2023qlora}. Qwen2.5-3B-Instruct est
chargé en 4 bits avec la quantification NF4, une double quantification et des
calculs en BF16. Les adaptateurs ont un rang \(r=16\), un facteur
\(\alpha=32\) et un dropout de 0,05. Ils ciblent les projections
\texttt{\footnotesize q\_proj}, \texttt{\footnotesize k\_proj}, \texttt{\footnotesize v\_proj} et
\texttt{\footnotesize o\_proj}, ainsi que \texttt{\footnotesize gate\_proj}, \texttt{\footnotesize up\_proj} et
\texttt{\footnotesize down\_proj} dans le bloc feed-forward.

L'entraînement utilise un lot unitaire et une accumulation de gradient pour
obtenir un lot effectif plus grand sans dépasser la mémoire disponible. Le
gradient checkpointing, l'optimiseur paginé et la désactivation du cache
complètent ce dispositif. Le jeu est séparé de manière déterministe entre
apprentissage et validation. La perte n'est calculée que sur les tokens produits
par l'assistant grâce aux masques du gabarit ; les messages système, demandes et
résultats d'outils servent de contexte, mais ne sont pas des cibles à recopier.


\subsection{Routage et état d'exécution sans session}

Le serveur expose cinq identifiants : un identifiant HMoE automatique et un
identifiant de diagnostic par expert. Le routage automatique inspecte la phase
et le rôle contenus dans la requête système. Le routage explicite reste utile
pour tester un adaptateur isolément, mais le tableau de bord ne doit connaître
que l'identifiant automatique. La spécialisation demeure ainsi interne au
service et ne complexifie pas la configuration d'une campagne.

Le modèle partage un état de poids, mais le suivi agentique est volontairement
sans session. À chaque requête, le serveur reconstruit l'avancement à partir de
la conversation OpenAI reçue : nombre d'appels réussis, tentatives de sauvegarde
rejetées, exigences encore manquantes et progression de l'intrusion. Il injecte
ensuite ce résumé autoritaire dans le dernier message. Aucune mémoire de campagne
n'est conservée entre deux requêtes, ce qui empêche le mélange de deux exécutions
concurrentes et facilite la reprise après redémarrage du service.

\subsection{Contrats d'outils, compactage et récupération}

La petite taille des experts rend les boucles d'outils sensibles au bruit. Pour
la reconnaissance, seuls cinq schémas sont rendus dans le prompt : découverte
ARP, découverte Nmap, scan Nmap, lecture et sauvegarde de livrable. Le registre
exécutable complet reste chez l'appelant ; la réduction concerne seulement ce
que voit le modèle. Dès que la couverture minimale est atteinte, le schéma visible
est réduit à la sauvegarde. À l'inverse, tout appel généré vers un outil absent
du contrat de la requête est supprimé avant d'être transmis au pipeline.

Le compactage est aligné sur les fenêtres d'entraînement. Le serveur remplace
les brouillons de livrables rejetés, résume les anciennes sorties réseau et, si
nécessaire, replie les vieux tours en un registre déterministe. Il préserve les
paires appel/résultat récentes et la dernière erreur contractuelle. La génération
est ensuite plafonnée selon l'expert et l'espace restant dans la fenêtre. Ce
mécanisme ne cherche donc pas seulement à réduire les tokens : il conserve les
informations dont dépend la prochaine action.

Lorsqu'une tentative de sauvegarde retourne une liste structurée
\texttt{\footnotesize missing\_requirements}, le serveur peut produire directement l'appel
manquant autorisé. Il transforme, par exemple, un port non couvert en scan Nmap
ciblé. Cette récupération déterministe intervient avant une nouvelle génération
et évite que le modèle répète le même livrable refusé. Dans les phases
d'intrusion, le pipeline maintient en plus un registre d'actions faisant autorité.
Si le modèle s'arrête ou si le fournisseur local devient indisponible, une
reprise bornée exécute uniquement des sondes déduites du contexte, sans inventer
d'identifiants. Le livrable final est alors reconstruit à partir du registre et
signale les couvertures encore incomplètes.

\subsection{Déploiement du service}

Le service HMoE charge le modèle de base quantifié et les quatre adaptateurs,
puis expose une API compatible OpenAI avec des routes de santé, de modèles et de
conversation. Il est lancé comme service utilisateur systemd. Sur la machine
GPU, l'écoute est limitée au pont Docker pour OpenWebUI ; l'accès depuis
\texttt{\footnotesize nato-master} passe par un proxy lié uniquement à l'adresse Tailscale.
Cette topologie évite d'exposer publiquement une API qui ne valide pas elle-même
un jeton Bearer. Après une période d'inactivité configurable, les allocations
CUDA inutilisées sont rendues au pilote sans interrompre le modèle chargé.


\section{Évolution détaillée du pipeline}

\subsection{Phase 2 : reconnaissance contractuelle}

La phase~2 transforme une cible réseau en surface d'attaque observée. L'agent
doit couvrir la découverte, les hôtes attendus et leurs services, puis enregistrer
\texttt{\footnotesize 02\_recon.md}. La sauvegarde n'est pas un simple signal de fin : l'outil
vérifie les appels effectués et renvoie les exigences manquantes. Les résultats
de scan sont projetés dans une représentation déterministe des hôtes, ports et
protocoles. Cette projection alimente la phase suivante même si la prose du
livrable est incomplète.

J'ai ajouté une réconciliation entre cette surface et le graphe de scénario.
Elle associe les observations aux nœuds connus et conserve les cibles non
résolues au lieu de les faire disparaître. En mode découverte, les liens peuvent
être inférés après la reconnaissance. Ainsi, la phase~3 ne dépend plus uniquement
d'une topologie préexistante ni de la capacité du modèle à recopier exactement
les résultats Nmap.

\subsection{Phase 3 : analyse par appareil}

La phase~3 est décomposée en trois étapes. Un scanner déterministe collecte
d'abord les indices disponibles pour chaque appareil. Des micro-agents spécialisés
analysent ensuite cette surface appareil par appareil, avec une limite de temps,
de tours et de tokens. Enfin, une agrégation déterministe fusionne les sorties
dans une file canonique de vulnérabilités. Une défaillance locale ne supprime pas
les résultats du scanner : elle est inscrite dans un artefact de statut et la
campagne continue avec un résultat partiel explicite.

L'analyse conserve les versions de logiciels et peut interroger un catalogue CVE
hors ligne. Les findings sont normalisés avec un identifiant stable, une cible,
un type, une sévérité, une preuve et un plan de vérification. Dans le profil HMoE
compact, certaines observations directes peuvent être retenues pour le rapport
sans être promues automatiquement vers l'exploitation. Ce cloisonnement évite
qu'une configuration seulement suspectée devienne une action offensive.

\subsection{Phase 4 : vérification par vulnérabilité}

La phase~4 planifie un travail indépendant pour chaque finding éligible. Le
contrôleur choisit un outil compatible avec le service et impose un contrat de
vérification. Le micro-agent ne reçoit donc pas une liberté d'exploitation
illimitée : il doit exécuter une sonde ciblée, rester dans le périmètre autorisé
et produire un résultat lié à cette sonde. Dans le profil compact, si l'appel au
modèle échoue ou si la sonde requise n'apparaît pas, le contrôleur peut exécuter
une sonde de repli bornée.

L'agrégateur distingue un succès, un échec conclusif, une erreur et l'absence de
test. Un succès enrichit la preuve ; un échec peut réfuter une simple suspicion ;
une erreur ne transforme pas la détection en faux positif. Les résultats sont
réunis dans \texttt{\footnotesize 04\_exploitation.json} avec leur provenance d'outil. Les
hôtes découverts durant cette phase déclenchent un cycle limité de reconnaissance
et d'analyse avant la phase~5, ce qui permet de traiter un pivot sans relancer
toute la campagne.

\subsection{Phase 5 : intrusion contrôlée}

La phase~5 raisonne sur les findings vérifiés, les chemins possibles et les
identifiants réellement récupérés. Un contexte structuré énumère les points
d'entrée, les cibles, les services et les chaînes candidates. Le contrat terminal
exige des actions observables : lecture du contexte, couverture des points
d'entrée, essais de réutilisation des identifiants sur les cibles concernées et,
lorsqu'un identifiant est disponible, au moins un résultat d'accès authentifié.

Chaque action est inscrite dans un registre faisant autorité. Le modèle produit
un mémo, mais il ne peut pas déclarer seul la phase complète. Le contrôleur
calcule la couverture, propose une action lorsqu'aucun progrès n'est observé et
limite les reprises à un nombre borné de tours. Après un accès authentifié, une
collecte post-accès récupère uniquement les éléments nécessaires à la preuve.
Une campagne qui n'exécute aucune action ou ne satisfait pas le contrat est
marquée bloquée ou incomplète ; elle n'est pas silencieusement convertie en
succès.

\subsection{Phase 6 : rapport fondé sur les artefacts}

La phase~6 ne réinjecte pas l'intégralité des traces. Le pipeline construit un
contexte compact à partir des vulnérabilités, résultats de vérification, chaînes
d'intrusion et observations de configuration. Les tables lourdes sont
pré-générées, et seuls les findings confirmés avec un niveau de preuve suffisant
sont présentés comme vérifiés. Le modèle se concentre ainsi sur l'analyse et la
hiérarchisation plutôt que sur la recopie de données volumineuses.

Le rapport généré est validé avant fusion avec les sections déterministes. Si le
modèle local échoue ou produit une synthèse incompatible avec les artefacts, un
repli compose le rapport à partir des données structurées. Cette dernière
frontière garantit qu'une formulation persuasive ne réintroduit pas une
vulnérabilité éliminée par la phase~4.

\section{Harness expérimental et politique d'évaluation \texttt{\footnotesize strict-v3}}

Le harness expérimental organise le déploiement, la collecte des artefacts et
la répétition des campagnes. Au sein de ce dispositif, l'évaluateur applique
la politique \texttt{\footnotesize strict-v3} pour mesurer non seulement la détection, mais aussi la
qualité de la correspondance et la traçabilité des preuves. Chaque scénario
associe une topologie, des packs de vulnérabilités et une vérité terrain. Des
contrats de matching revus décrivent les champs structurants et les
compatibilités autorisées. Le moteur construit les correspondances candidates,
puis résout une association bipartite globale : un finding ne peut expliquer
qu'une entrée de vérité terrain, et réciproquement. Cette règle supprime les
scores artificiels où une même observation créditerait plusieurs vulnérabilités.

Dans cette nomenclature, \texttt{\footnotesize strict-v3} désigne la politique de scoring,
\texttt{\footnotesize strict-v3.3} le contrat métrique, \texttt{\footnotesize strict-v3.2} le schéma des
contrats de matching et \texttt{\footnotesize evidence-v2} le contrat de preuve. Ces
versions séparées permettent de faire évoluer un contrat sans changer
implicitement le sens des autres résultats.

Le crédit dépend de la précision structurelle du match. Une correspondance
exacte reçoit le crédit maximal ; un type correct mais incomplet ou une
compatibilité explicitement prévue reçoit un crédit moindre. Le score ajusté à
la qualité combine ce crédit avec l'erreur de sévérité et la qualité de
vérification. Le \emph{verified F1} impose en outre un appel d'outil lié au
finding et dont le résultat est explicitement positif. Les métriques de chemins
exigent que toutes les vulnérabilités attendues soient retrouvées ; leurs
variantes vérifiées contrôlent aussi l'ordre des appareils dans une chaîne de
phase~5.

J'ai appliqué une politique d'échec fermé aux métadonnées et aux artefacts. Une
campagne incomplète reste incomplète. Une version de contrat métrique incompatible
neutralise les scores dépendant des preuves au lieu de les calculer avec une
sémantique différente. Pour le split scellé, le worker ne reçoit ni vérité
terrain ni détails privés ; le contrôleur renvoie seulement un résumé signé. Les
scénarios attendus mais absents valent zéro lors de l'agrégation officielle.

L'agrégation évite enfin qu'un scénario riche en findings domine le résultat.
Les répétitions sont d'abord moyennées pour chaque scénario, puis les scénarios
sont macro-moyennés à poids égal. Les métriques brutes, créditées,
ajustées et vérifiées restent séparées afin qu'une amélioration de précision ne
masque pas une régression de preuve. Un juge sémantique peut aider au diagnostic,
mais la correspondance un-à-un, la validation de schéma et les scores finaux
restent déterministes, ce qui limite les biais du paradigme
\emph{LLM-as-a-judge} \cite{zheng2023judge}.

\section{Maîtrise du contexte}

J'ai borné les sorties SSH, compacté les résumés de topologie et dédoublonné
certaines observations MQTT ou HTTP. Ces optimisations ont été appliquées avec
prudence : une première réduction trop agressive de la surface d'attaque a été
annulée parce qu'elle supprimait une information utile. Ce retour a directement
guidé le compactage du serveur HMoE, qui conserve désormais les contrats, les
erreurs récentes et le registre des outils avant de réduire la prose.

\paragraph{Exploration topologique initiale.}
Au début du stage, j'ai également relié au pipeline un graphe d'attaque orienté
pondéré. Le coût d'une arête est l'inverse d'un poids heuristique de
franchissement. Après une découverte ou une compromission, les distances vers la
cible sont recalculées ; leur variation localise les nœuds rapprochés de
l'objectif. Une propagation sur les prédécesseurs forme un cône d'impact, dans
l'esprit des graphes d'attaque multi-hôtes \cite{ou2005mulval}. Ces valeurs ne
sont pas des probabilités calibrées : elles servent uniquement à prioriser une
analyse. Cette expérimentation est restée un apport secondaire et ne constitue
pas l'axe principal des campagnes présentées dans ce rapport.

\section{Résultats et difficultés}
\subsection{Instantané expérimental du 24 août 2026}
\label{sec:resultats-provisoires}

Les résultats présentés ici constituent un instantané de travail et non une
campagne finale. J'ai retenu quatre exécutions terminées et comparables au
regard des métadonnées disponibles, accessibles sur \texttt{nato-master}.
Elles utilisent toutes le modèle MiniMax-M2.7, le profil \texttt{full}, le
split \texttt{dev-public}, la révision logicielle \texttt{e7c4f6a} et la
politique d'évaluation \texttt{strict-v3} avec le contrat métrique
\texttt{strict-v3.3}. Les métadonnées indiquent que le modèle n'avait pas accès
à la vérité terrain. Chaque scénario ne comporte qu'une répétition ; les
valeurs ne permettent donc pas encore d'estimer la variance.

\begin{table}[H]
  \centering
  \small
  \caption{Résultats provisoires des quatre campagnes sélectionnées}
  \label{tab:resultats-provisoires}
  \begin{tabular}{lrrrrr}
    \toprule
    Scénario & Précision & Rappel & F1 & F1 vérifié & F1 qualité \\
    \midrule
    S1 & \(0{,}875\) & \(0{,}583\) & \(0{,}700\) & \(0{,}600\) & \(0{,}625\) \\
    S2 & \(0{,}833\) & \(0{,}769\) & \(0{,}800\) & \(0{,}640\) & \(0{,}670\) \\
    S5 & \(0{,}900\) & \(0{,}600\) & \(0{,}720\) & \(0{,}640\) & \(0{,}680\) \\
    S8 & \(0{,}533\) & \(0{,}571\) & \(0{,}552\) & \(0{,}414\) & \(0{,}448\) \\
    \midrule
    Moyenne macro & \(0{,}785\) & \(0{,}631\) & \(0{,}693\) & \(0{,}574\) & \(0{,}606\) \\
    \bottomrule
  \end{tabular}
\end{table}

Comme l'indique \cref{tab:resultats-provisoires}, la moyenne macro atteint un
F1 de \(0{,}693\), contre \(0{,}574\) lorsque seuls les constats vérifiés sont
crédités. Le F1 ajusté par la qualité s'établit à \(0{,}606\). L'écart entre
ces mesures suggère qu'une détection correcte au niveau du type et de la cible
ne s'accompagne pas toujours d'une preuve d'exécution suffisamment forte.

La complétion moyenne de la phase 4 atteint \(92{,}8\,\%\) et la couverture
d'exploitation \(82{,}4\,\%\). En revanche, la couverture macro des chemins
n'est que de \(52{,}1\,\%\), et celle des chemins entièrement vérifiés de
\(6{,}3\,\%\). S1 est le seul scénario de cet échantillon à obtenir une
couverture vérifiée non nulle, égale à \(25\,\%\). Dans cet échantillon, la
difficulté principale semble donc être la conservation d'une chaîne de preuves
complète sur plusieurs étapes, et non le seul lancement des vérifications
locales.

S2 obtient le meilleur F1 de l'échantillon avec \(0{,}800\). S8 est le cas le
plus difficile : huit vrais positifs, sept faux positifs et six faux négatifs
conduisent à un F1 de \(0{,}552\). Ce résultat motive les gardes consacrées à
la précision sur les scénarios étendus et à la réduction des hallucinations.

Les quatre exécutions totalisent environ 3,58 dollars, 11,0 millions de tokens,
1\,202 appels d'outils et 30 erreurs d'outils. Elles utilisent cependant un
modèle généraliste distant et non le \gls{hmoe} local. Cet instantané constitue
donc un test d'intégration du pipeline et du harness ; il ne démontre
pas encore l'efficacité des experts entraînés. Une campagne finale devra figer
la version du code, répéter chaque scénario et publier moyenne, dispersion et
intervalle de confiance pour les baselines et le \gls{hmoe}.

Les identifiants de runs utilisés sont
\texttt{2026-08-24\_092203} pour S1,
\texttt{2026-08-21\_144141} pour S2,
\texttt{2026-08-24\_083553} pour S5 et
\texttt{2026-08-21\_151305} pour S8.


\subsection{Difficultés techniques rencontrées}

\begin{description}
  \item[Aligner apprentissage et exécution] Un adaptateur entraîné sur un outil,
  un format ou une fenêtre différents de la production peut échouer malgré une
  perte d'entraînement correcte. J'ai donc partagé le gabarit, les schémas et les
  budgets entre préparation, préflight et service.

  \item[Préserver la preuve sous contrainte de contexte] Une réduction aveugle
  économise des tokens mais peut séparer une conclusion de son appel d'outil. Le
  découpage atomique de la phase~5 et le registre déterministe répondent à ce
  risque.


  \item[Distinguer terminaison et succès] Un livrable sauvegardé ou une réponse
  fluide ne prouve pas la couverture. Les contrats de phase, les reprises bornées
  et l'échec fermé rendent les lacunes visibles au lieu de les masquer.

\end{description}
